% =====================================================================================
% QUESTIONNAIRE D'ELICITATION -- version 2, 13 aout 2026
%
% CE QUI A CHANGE PAR RAPPORT A LA VERSION DE JUILLET, ET POURQUOI. Le paquet est destine a
% etre soumis a un actuaire de l'Institut avant tout envoi. Six defauts le rendaient
% refusable en l'etat :
%
%  1. LA GRAINE SUR xi ETAIT FAUSSE, et c'est le defaut le plus grave. Elle annoncait 0,90
%     comme valeur vraie, or 0,90 est la valeur POSEE du pire cas ; la calibration figee du
%     memoire donne xi = 0,5954, avec un IC90 de [0,30 ; 0,83]. Dans la methode de Cooke une
%     graine fausse ne degrade pas le resultat a la marge : elle INVERSE les poids, un expert
%     juste etant penalise. Et meme corrigee a 0,5954 la question resterait invalide comme
%     graine, une graine devant avoir une valeur REALISEE et non une estimation assortie d'un
%     intervalle deux fois plus large que la valeur elle-meme. Elle est donc retiree de la
%     partie A. Les neuf autres graines ont ete verifiees une par une contre les sorties de
%     scripts versionnees : toutes justes.
%  2. LA CONVENTION DE DIRECTION ETAIT INVERSEE par rapport au memoire. Le formulaire disait
%     « lien de k vers j » quand le memoire ecrit W_jk avec j la SOURCE et k la CIBLE (la
%     ligne emet, la colonne recoit). Une saisie fidele au formulaire aurait TRANSPOSE W.
%     Aligne : partout j -> k, j emet.
%  3. Aucune mention d'usage, de conservation ni d'anonymat, ce qu'aucun repondant externe
%     n'accepte aujourd'hui.
%  4. Aucune logistique : ni version, ni retour, ni delai, ni contact.
%  5. Aucun exemple traite, alors que le malentendu sur les quantiles est le premier facteur
%     de mauvaises reponses.
%  6. La graine sur xi portait un indice entre parentheses qui soufflait l'ordre de grandeur.
%     Le retrait de la question regle aussi ce point.
% =====================================================================================
\documentclass[11pt,a4paper]{article}

\usepackage[utf8]{inputenc}
\usepackage[T1]{fontenc}
\usepackage{lmodern}
\usepackage[french]{babel}
\usepackage[margin=1.7cm]{geometry}
\usepackage{microtype}
\usepackage{array}
\usepackage{booktabs}
\usepackage[table]{xcolor}
\usepackage[most]{tcolorbox}
\usepackage{enumitem}
\usepackage{fancyhdr}
% longtable et non tabular : la table de la partie A compte neuf lignes et un tabular, qui est
% insecable, ne tenait pas dans le reste de la premiere page. Il partait donc en entier sur la
% suivante en laissant une demi-page blanche, et le formulaire passait a trois pages.
\usepackage{longtable}

\definecolor{navy}{HTML}{1F3864}
\definecolor{lightblue}{HTML}{D6E0F0}
\definecolor{accent}{HTML}{C0491F}
\definecolor{softgray}{HTML}{F2F2F2}

\newtcolorbox{cle}[1][]{colback=lightblue!40,colframe=navy,boxrule=0.6pt,
  arc=2pt,left=8pt,right=8pt,top=6pt,bottom=6pt,#1}
\newtcolorbox{mention}[1][]{colback=softgray,colframe=black!35,boxrule=0.5pt,
  arc=2pt,left=8pt,right=8pt,top=5pt,bottom=5pt,#1}

% 1,75 et non 2,0 : a 2,0 la table de la partie A ne tenait plus sur la premiere page et y
% laissait une demi-page blanche, le formulaire passant a trois pages. La hauteur de cellule
% reste suffisante pour une reponse manuscrite.
\renewcommand{\arraystretch}{1.55}
\newcolumntype{A}{>{\centering\arraybackslash}m{1.7cm}}

\pagestyle{fancy}
\fancyhf{}
\lfoot{\footnotesize Élicitation $W$ --- version 2.0 du 13 août 2026}
\rfoot{\footnotesize page \thepage/\pageref{LastPage}}
\renewcommand{\headrulewidth}{0pt}
\usepackage{lastpage}

\title{\vspace{-1.3cm}\color{navy}\bfseries Questionnaire d'élicitation de $W$\\[2pt]
  \large Méthode de Cooke --- une séance d'environ 30 minutes}
\date{}

\begin{document}
\maketitle
\thispagestyle{fancy}   % sinon \maketitle repasse la premiere page en style plain
\vspace{-1.4cm}

\noindent\textbf{Expert :} \underline{\hspace{4.2cm}} \quad
\textbf{Fonction :} \underline{\hspace{4.2cm}} \quad
\textbf{Date :} \underline{\hspace{2.6cm}}

\vspace{0.25cm}
\begin{cle}
\textbf{Comment répondre.} Pour chaque quantité, donnez \textbf{trois nombres} : une valeur
\textbf{basse (5\,\%)}, une valeur \textbf{médiane (50\,\%)} et une valeur \textbf{haute
(95\,\%)}, telles que vous estimiez à \textbf{$90\,\%$ de chances} que la vraie valeur tombe
entre votre $5\,\%$ et votre $95\,\%$. On mesure votre jugement, pas vos connaissances, et il
n'y a pas de réponse attendue. \textbf{Répondez même très incertain} : un intervalle large est
une réponse honnête et exploitable, une case vide ne l'est pas. Écrivez la \textbf{haute
d'abord}, puis la \textbf{basse}, et la \textbf{médiane en dernier} : commencer par la médiane
ancre les deux autres autour d'elle et produit des intervalles trop étroits, seul défaut que la
méthode sanctionne réellement.
\end{cle}

\vspace{0.2cm}
\begin{mention}
\footnotesize
\textbf{Exemple traité, sur une question sans rapport avec le sujet.} \emph{Quelle est la
hauteur de la tour Montparnasse, en mètres ?} Une réponse possible : haute $=260$
(\og je serais surpris, une chance sur vingt, que ce soit plus haut \fg), basse $=150$
(\og de même en dessous \fg), médiane $=200$ (\og autant de chances au-dessus qu'en
dessous \fg). L'intervalle $[150\,;260]$ est large, et c'est bien ainsi : il annonce
franchement ce que l'on ignore.
\end{mention}

\section*{Partie A --- Questions d'étalonnage}
\footnotesize\emph{Neuf quantités effectivement mesurées dans les données de l'étude. Leur
valeur est connue de l'analyste et n'est pas affichée ici ; chacune est reproductible par un
script versionné, de sorte qu'aucune notation ne repose sur un chiffre invérifiable. Ces
questions servent à \textbf{pondérer} les experts, non à les départager sur leurs
connaissances : le périmètre est celui d'une base publique d'incidents du secteur financier,
et personne n'est censé en connaître les valeurs de mémoire.}
\normalsize

\vspace{0.2cm}
\begin{longtable}{|p{9.6cm}|A|A|A|}
\hline
\rowcolor{lightblue!50}
\textbf{Quantité} & \textbf{5\,\%} & \textbf{50\,\%} & \textbf{95\,\%} \\ \hline
\endfirsthead
\hline
\rowcolor{lightblue!50}
\textbf{Quantité \emph{(suite)}} & \textbf{5\,\%} & \textbf{50\,\%} & \textbf{95\,\%} \\ \hline
\endhead
A1. Part des sinistres financiers portés par une cause commune, c'est-à-dire imputables à un
même prestataire tiers, en \% & & & \\ \hline
A2. Délai médian entre la survenance d'un incident et sa déclaration, en jours & & & \\ \hline
A3. Part des incidents déclarés plus d'un mois après leur survenance, en \% & & & \\ \hline
A4. Durée médiane de \emph{containment}, de la brèche à sa fermeture, en jours & & & \\ \hline
A5. Nombre d'incidents TIC matériels par an, pour une entité financière \textbf{de grande
taille} (au moins dix incidents sur la période) & & & \\ \hline
A6. Délai médian de déclaration pour les seuls incidents de type \emph{piratage}, en jours & & & \\ \hline
A7. Indice de Gini de la concentration des sinistres sur les jours de l'année
($0$ = un même nombre chaque jour ; $1$ = tout un seul jour) & & & \\ \hline
A8. Nombre maximal d'entités victimes d'une même cause commune en une seule journée & & & \\ \hline
A9. Part des jours de l'année qui portent une cause commune, en \% & & & \\ \hline
\end{longtable}

\section*{Partie B --- Liens dirigés entre piliers}
\footnotesize\emph{\textbf{Force du lien dirigé de $j$ vers $k$}, notée $W_{jk}$, entre
$\mathbf{0}$ et $\mathbf{1}$. La \textbf{première} lettre est celle du pilier qui
\textbf{défaille en premier}, la seconde celle du pilier \textbf{entraîné}. Lecture :
$W_{jk}=0$ signifie qu'une défaillance de $j$ n'entraîne \emph{jamais} celle de $k$ ;
$W_{jk}=1$ qu'elle l'entraîne \emph{quasi certainement} dans la foulée. Donnez là aussi
$5\,\% / 50\,\% / 95\,\%$. \textbf{Les deux sens d'une même paire sont posés séparément et
n'ont aucune raison d'être égaux} : c'est précisément l'asymétrie que l'étude cherche à
mesurer, donc ne cherchez pas à les faire correspondre.}\\[3pt]
\textbf{Les cinq piliers :} P1 gouvernance et gestion du risque TIC \textbullet{} P2 gestion
des incidents et réponse \textbullet{} P3 tests de résilience et détection \textbullet{}
P4 risque lié aux prestataires tiers \textbullet{} P5 partage d'information.
\normalsize

\vspace{0.2cm}
\begin{longtable}{|p{9.6cm}|A|A|A|}
\hline
\rowcolor{lightblue!50}
\textbf{Lien dirigé $j \to k$, force entre $0$ et $1$} & \textbf{5\,\%} & \textbf{50\,\%} & \textbf{95\,\%} \\ \hline
\endfirsthead
\hline
\rowcolor{lightblue!50}
\textbf{Lien dirigé $j \to k$ \emph{(suite)}} & \textbf{5\,\%} & \textbf{50\,\%} & \textbf{95\,\%} \\ \hline
\endhead
B1. \textbf{P1 $\to$ P2} : une défaillance de gouvernance entraîne un incident mal géré & & & \\ \hline
B2. \textbf{P2 $\to$ P1} : \emph{le sens inverse du précédent} --- un incident mal géré entraîne une défaillance de gouvernance & & & \\ \hline
B3. \textbf{P4 $\to$ P2} : la défaillance d'un prestataire tiers entraîne un incident mal géré & & & \\ \hline
B4. \textbf{P1 $\to$ P4} : une défaillance de gouvernance entraîne une défaillance côté prestataires & & & \\ \hline
B5. \textbf{P3 $\to$ P1} : un défaut de tests et de détection entraîne une défaillance de gouvernance & & & \\ \hline
B6. \textbf{P5 $\to$ P2} : un défaut de partage d'information entraîne un incident mal géré & & & \\ \hline
\end{longtable}

\vspace{0.05cm}
\footnotesize\emph{Le modèle compte $20$ liens dirigés ; six seulement sont élicités ici pour
tenir la séance en trente minutes, à savoir les quatre que le corpus documentaire atteste et la
paire P1\,/\,P2 dans ses deux sens, qui porte la question d'asymétrie. Les quatorze autres
restent bornés par l'identification partielle, et les éliciter serait une extension du même
formulaire, sans changement de méthode.}
\normalsize

\vspace{0.18cm}
\begin{cle}
\textbf{Ce qui est fait de vos réponses.} La partie A mesure votre \emph{calibration} (vos
intervalles contiennent-ils la vraie valeur à la fréquence annoncée) et votre \emph{information}
(leur finesse) ; ces deux scores fixent votre poids. La partie B, pondérée par ce poids et
croisée avec les autres répondants, produit chaque lien de $W$ sous forme d'\textbf{intervalle},
jamais de point. Le traitement est automatique (\texttt{26\_elicitation\_cooke.py},
\texttt{26b\_elicitation\_reelle.py}). Aucune donnée d'incident propre à votre organisation
n'est demandée, uniquement votre jugement.
\end{cle}

\vspace{0.18cm}
\begin{mention}
\footnotesize
\textbf{Usage, conservation et anonymat.} Ce questionnaire s'inscrit dans un mémoire d'actuariat
(ENSAE Paris, Institut des Actuaires, promotion 2026) sur la quantification du besoin de capital
lié à la non-conformité au règlement DORA, et vos réponses servent à cette seule fin. Elles sont
enregistrées sous un \textbf{identifiant anonyme} : ni votre nom ni celui de votre organisation
n'apparaissent dans le mémoire, dans les fichiers versionnés, ni dans aucune restitution ; seuls
les intervalles agrégés sont publiés, et les poids individuels sous cet identifiant. Aucune
donnée personnelle autre que votre fonction n'est collectée. Vous pouvez demander le retrait de
vos réponses jusqu'à la remise du mémoire, sans avoir à le motiver : elles seront alors détruites
et exclues de l'agrégation. Conservation jusqu'à la soutenance, puis destruction.
\textbf{Retour} à \underline{\hspace{4.6cm}} avant le \underline{\hspace{2.4cm}}. Une séance
accompagnée de trente minutes est proposée et améliore sensiblement la qualité des intervalles,
mais un envoi rempli seul reste exploitable. Le protocole détaillé, même version, est joint.
\end{mention}

\vspace{0.45cm}
\section*{Accord du répondant}
\vspace{-0.15cm}

\noindent\framebox[0.42cm]{\rule{0pt}{0.32cm}}\; J'accepte que mes réponses soient utilisées,
sous identifiant anonyme, dans le cadre du mémoire décrit ci-dessus.

\vspace{0.18cm}
\noindent\framebox[0.42cm]{\rule{0pt}{0.32cm}}\; Je souhaite recevoir la synthèse des résultats
agrégés une fois l'agrégation faite.

\vspace{0.18cm}
\noindent\framebox[0.42cm]{\rule{0pt}{0.32cm}}\; Je souhaite relire la façon dont mes réponses
sont citées avant toute diffusion.

\vspace{0.5cm}
\noindent\textbf{Fait à} \underline{\hspace{4.4cm}} \quad
\textbf{le} \underline{\hspace{3.2cm}} \quad
\textbf{Signature} \underline{\hspace{4.4cm}}

\vspace{0.7cm}
\noindent\rule{\linewidth}{0.4pt}

\vspace{0.2cm}
\footnotesize\emph{Réservé à l'analyste, à ne pas remplir par le répondant.} Identifiant
attribué : \underline{\hspace{2.4cm}} \quad Date de réception : \underline{\hspace{2.4cm}}
\quad Saisi dans \texttt{elicitation\_reponses.csv} : \framebox[0.34cm]{\rule{0pt}{0.26cm}}
\quad Séance accompagnée : \framebox[0.34cm]{\rule{0pt}{0.26cm}}
\normalsize

\end{document}
